% from ega4-18-fr.tex
\oldpage[IV-4]{208}

d'autre part, $(Q_k')_{i+1}$ est l'image de $\sum_{h=0}^{k-i-1}\mathfrak J^{k-i-h-1}M_{i+h+1}$ dans $M_k'/M_{k+1}'$. Écrire que $(\psi_M)_{ki}$ est injectif équivaut donc à écrire la condition $(\mathrm I_{k-i,i})$; d'où notre assertion.
\end{env}

\begin{env}[19.5.5.2]
\label{IV.19.5.5.2}
Pour que la suite $\mathbf f$ soit $\operatorname{gr}_\bullet(M)$-quasi-régulière, il suffit que, pour $p\geq0$ et $q\geq0$, la condition suivante soit vérifiée:

$(\mathrm S_{q,p})$ Pour toute famille $(x_{\mathbf q}^p)_{|\mathbf q|=q}$ d'éléments de $M_p$, la relation
\[
  \sum_{|\mathbf q|=q}\mathbf f^{\mathbf q}x_{\mathbf q}^p\in M_{p+1}
\]
implique $x_{\mathbf q}^p\in M_{p+1}+\mathfrak JM_p$ pour tout $\mathbf q$ tel que $|\mathbf q|=q$.

Par définition, dire que $\mathbf f$ est $\operatorname{gr}_\bullet(M)$-quasi-régulière signifie que, pour tout $p\geq0$ et tout $q\geq0$, une relation
\[
  \sum_{|\mathbf q|=q}\mathbf f^{\mathbf q}x_{\mathbf q}^p\in M_{p+1}+\mathfrak J^{q+1}M_p
\]
pour une famille $(x_{\mathbf q}^p)_{|\mathbf q|=q}$ d'éléments de $M_p$ entraîne $x_{\mathbf q}^p\in M_{p+1}+\mathfrak JM_p$ pour tout $\mathbf q$ tel que $|\mathbf q|=q$. Or l'hypothèse sur la suite $(x_{\mathbf q}^p)$ signifie qu'il existe une famille $(y_{\mathbf q}^p)_{|\mathbf q|=q}$ d'éléments de $\mathfrak JM_p$ telle que l'on ait $\sum_{|\mathbf q|=q}\mathbf f^{\mathbf q}(x_{\mathbf q}^p-y_{\mathbf q}^p)\in M_{p+1}$. La condition $(\mathrm S_{q,p})$ entraîne alors $x_{\mathbf q}^p-y_{\mathbf q}^p\in M_{p+1}+\mathfrak JM_p$ pour tout $\mathbf q$ tel que $|\mathbf q|=q$, d'où $x_{\mathbf q}^p\in M_{p+1}+\mathfrak JM_p$; on voit donc que les conditions $(\mathrm S_{q,p})$ entraînent que la suite $\mathbf f$ est $\operatorname{gr}_\bullet(M)$-quasi-régulière.
\end{env}

Il suffit donc de prouver la proposition suivante:

\begin{env}[19.5.5.3]
\label{IV.19.5.5.3}
Pour tout $r\geq0$, si les conditions $(\mathrm I_{q,p})$ sont vérifiées pour $p+q<r$, alors les conditions $(\mathrm S_{q,p})$ sont vérifiées pour $p+q<r$.

Introduisons les conditions

$(\mathrm A_{q,p})$ Pour toute famille $(x_{\mathbf q}^p)_{|\mathbf q|=q}$ d'éléments de $M_p$, la relation
\[
  \sum_{|\mathbf q|=q}\mathbf f^{\mathbf q}x_{\mathbf q}^p\in M_{p+1}
\]
implique
\[
  \sum_{|\mathbf q|=q}\mathbf f^{\mathbf q}x_{\mathbf q}^p\in\sum_{i=1}^q\mathfrak J^{q-i}M_{p+i}.
\]

Il est clair que la conjonction de $(\mathrm A_{q,p})$ et de $(\mathrm I_{q,p})$ implique $(\mathrm S_{q,p})$. D'autre part, les conditions $(\mathrm A_{1,p})$ et $(\mathrm S_{0,p})$ sont triviales. Par suite, (19.5.5.3) résultera de la proposition suivante:
\end{env}

\begin{env}[19.5.5.4]
\label{IV.19.5.5.4}
Quels que soient $p\geq0$ et $q\geq1$, la condition $(\mathrm A_{q+1,p})$ est entraînée par la conjonction des conditions $(\mathrm A_{q,p})$, $(\mathrm I_{q-1,p+1})$, $(\mathrm I_{q-2,p+2})$, $\ldots$, $(\mathrm I_{0,p+q})$.

En effet, une fois prouvée cette proposition, les conditions $(\mathrm I_{q,p})$ supposées vérifiées pour $p+q<r$ entraîneront, pour chaque $p\leq r$, la condition $(\mathrm A_{q,p})$ pour $1\leq q<r-p$, par récurrence sur $q$.

Prouvons donc (19.5.5.4). Notons que la condition $(\mathrm A_{q,p})$ équivaut aussi à
\end{env}

\begin{env}[19.5.5.5]
\label{IV.19.5.5.5}
\[
  \mathfrak J^qM_p\cap M_{p+1}\subset\sum_{i=1}^q\mathfrak J^{q-i}M_{p+i}.
\]

Soit donc $m\in\mathfrak J^{q+1}M_p\cap M_{p+1}\subset\mathfrak J^qM_p\cap M_{p+1}$. Appliquant la condition $(\mathrm A_{q,p})$, on voit qu'il existe, pour chaque $i$ tel que $1\leq i\leq q$, une famille $(y_{\mathbf q}^{p+i})_{|\mathbf q|=q-i}$ d'éléments de $M_{p+i}$ telle que
\[
  m=\sum_{i=1}^q\left(\sum_{|\mathbf q|=q-i}\mathbf f^{\mathbf q}y_{\mathbf q}^{p+i}\right).
\]
Supposons que, pour $1\leq i<j$ ($j\geq1$), on ait prouvé que
\[
  y_{\mathbf q}^{p+i}\in\mathfrak JM_{p+i}+M_{p+i+1}\qquad(\text{pour }|\mathbf q|=q-i).
\]
On en déduit par définition (19.5.4)
\[
  m-\sum_{i<j}\left(\sum_{|\mathbf q|=q-i}\mathbf f^{\mathbf q}y_{\mathbf q}^{p+i}\right)\in M_{p+q+1}'.
\]
Puisque par hypothèse $(\mathrm I_{q-j,p+j})$ est vérifiée, on en déduit
\[
  y_{\mathbf q}^{p+j}\in\mathfrak JM_{p+j}+M_{p+j+1}
\]

